\documentclass[11pt]{article}

\usepackage[utf8]{inputenc}
\usepackage[T1]{fontenc}
\usepackage[ngerman]{babel}
\usepackage{amsmath, amssymb, amsthm, amsfonts}
\usepackage[hidelinks]{hyperref}
\usepackage{geometry}
\geometry{a4paper, margin=1in}

\title{Der mathematische Beweis der Periodenformel\\ \large Schritt-für-Schritt erklärt für Ingenieure}
\author{Dirk Kunert \\ \small(Fokus auf das mathematische Verständnis des Beweises)}
\date{\today}

\begin{document}

\maketitle

\begin{abstract}
Dieses Dokument ist ein tieferer Einstieg in den eigentlichen mathematischen Beweis der Periodenlänge $N$ (für den Fall $N < D$). Es führt durch die fünf Beweisschritte des Original-Papers, erklärt jedoch die verwendeten abstrakten algebraischen Konzepte (wie Modulo-Arithmetik, Teilerfremdheit und modulare Inverse) anschaulich an physikalisch-technischen Analogien wie Zahnrädern und linearen Gleichungssystemen.
\end{abstract}

\section{Wichtige mathematische Werkzeuge (Crashkurs)}

Bevor wir in den Beweis einsteigen, müssen drei Fachbegriffe der Zahlentheorie geklärt werden, auf denen der Beweis aufbaut:

\begin{enumerate}
    \item \textbf{Modulo-Arithmetik ($x \equiv y \pmod D$):}\\
    Das ist die "Uhrzeiger-Mathematik". Zwei Zahlen sind "gleich modulo $D$", wenn sie bei der Division durch $D$ denselben Rest lassen. \\
    \textit{Ingenieurs-Analogie:} Stellen Sie sich ein Zahnrad mit $D$ Zähnen vor (nummeriert von $0$ bis $D-1$). Wenn Sie das Rad um $x$ Zähne weiterdrehen, landen Sie auf demselben Zahn, als wenn Sie es um $y$ Zähne drehen. Das bedeutet $x$ und $y$ sind modulo $D$ äquivalent.
    
    \item \textbf{Teilerfremdheit ($\gcd(a, b) = 1$):}\\
    Zwei Zahlen sind teilerfremd (größter gemeinsamer Teiler ist $1$), wenn sie sich nicht kürzen lassen. \\
    \textit{Ingenieurs-Analogie:} Wenn ein Ritzel mit $a$ Zähnen ein Rad mit $b$ Zähnen antreibt und $\gcd(a,b)=1$ gilt, dann berührt im Laufe der Zeit \textit{jeder} Zahn des Ritzels \textit{jeden} Zahn des Rades. Es gibt keinen lokalen Verschleiß.
    
    \item \textbf{Das modulare Inverse ($a^{-1} \pmod D$):}\\
    In den reellen Zahlen ist das Inverse zu $3$ einfach $1/3$, denn $3 \cdot (1/3) = 1$. In der Modulo-Arithmetik (wo es nur ganze Zahlen gibt), ist das Inverse zu $a$ eine ganze Zahl $x$, sodass das Produkt $a \cdot x$ bei Division durch $D$ den Rest $1$ lässt: $a \cdot x \equiv 1 \pmod D$. Eine solche Zahl existiert \textbf{nur dann}, wenn $a$ und $D$ teilerfremd sind ($\gcd(a, D) = 1$).
    
    \begin{itemize}
        \item \textit{Beispiel (Inverse existiert):} Wir suchen $3^{-1} \pmod 7$. Wir probieren die möglichen Werte für $x$: \\
        $3 \cdot 1 = 3 \equiv 3$, \quad $3 \cdot 2 = 6 \equiv 6$, \quad $3 \cdot 3 = 9 \equiv 2$, \\
        $3 \cdot 4 = 12 \equiv 5$, \quad \textbf{$3 \cdot 5 = 15 \equiv 1 \pmod 7$}. \\
        Da $3 \cdot 5$ den Rest $1$ lässt, ist $5$ das modulare Inverse von $3$ modulo $7$.
        
        \item \textit{Beispiel (Inverse existiert nicht):} Wir suchen $2^{-1} \pmod 4$. Wir probieren: \\
        $2 \cdot 1 = 2 \equiv 2$, \quad $2 \cdot 2 = 4 \equiv 0$, \quad $2 \cdot 3 = 6 \equiv 2$, \quad $2 \cdot 4 = 8 \equiv 0 \pmod 4$. \\
        Wir erhalten immer gerade Reste. Da $\gcd(2, 4) = 2 \neq 1$, lässt sich der Rest $1$ niemals erreichen. Das Inverse existiert hier nicht.

        \item \textit{Ingenieurs-Analogie:} Wenn ein Ritzel mit $a$ Zähnen ein Rad mit $D$ Zähnen antreibt und $\gcd(a,D)=1$ ist, dann gibt es genau eine bestimmte Anzahl von Umdrehungen $x$ des Ritzels, nach der das große Rad \textit{genau einen Zahn} weiter gedreht ist als bei seiner Ausgangsposition. Diese Anzahl an Umdrehungen $x$ entspricht dem modularen Inversen!
    \end{itemize}
\end{enumerate}

\section{Schritt 1: Die Invariante $r$ (Die "Spuren")}

Unser Toleranzband wird durch die Ungleichung definiert:
\[
    -\beta\omega \;\le\; \alpha x - \beta y \;\le\; \alpha\omega
\]
Der Ausdruck $r := \alpha x - \beta y$ ist extrem wichtig. Er ist eine Konstante (eine Invariante) für alle Punkte, die auf einer parallelen Linie zur Hauptachse liegen. Die Hauptachse verläuft in Richtung des Vektors $(\beta, \alpha)$. Wenn man von einem Gitterpunkt $(x,y)$ genau in diese Richtung zum nächsten Punkt $(x+\beta, y+\alpha)$ geht, ändert sich $r$ nicht:
\[
    \alpha(x+\beta) - \beta(y+\alpha) = \alpha x + \alpha\beta - \beta y - \beta\alpha = \alpha x - \beta y = r
\]
Jeder ganzzahlige Wert von $r$, der in das Intervall $[-\beta\omega, \alpha\omega]$ fällt, repräsentiert also eine kontinuierliche "Spur" (Track) von Gitterpunkten. Die Anzahl der ganzen Zahlen in diesem Intervall ist genau unsere Konstante $N$. Es gibt also $N$ Spuren.

\section{Schritt 2: Das lineare Gleichungssystem}

Für jede Spur (festes $r$) wollen wir wissen, an welchen Orten $\tilde{p}$ die Gitterpunkte auf unserer Projektionsachse landen. Der Ort $\tilde{p}$ ist durch die Projektion definiert als $\tilde{p} = \beta x + \alpha y$. 
Das ergibt ein Gleichungssystem mit zwei Variablen ($x, y$):
\begin{align*}
    \beta x + \alpha y &= \tilde{p} \quad (\text{Position auf der Achse})\\
    \alpha x - \beta y &= r \quad (\text{Nummer der Spur})
\end{align*}
Das lösen wir ganz klassisch nach $x$ und $y$ auf. Die Determinante ist $-\alpha^2 - \beta^2 = -D$. Das ergibt:
\[
    x = \frac{\beta\tilde{p} + \alpha r}{D}, \qquad y = \frac{\alpha\tilde{p} - \beta r}{D}
\]
Jetzt kommt der Trick: Die Punkte $(x,y)$ müssen auf dem Gitter liegen, also \textbf{ganze Zahlen} sein! Das bedeutet, der Zähler muss durch den Nenner ($D$) ohne Rest teilbar sein. In Modulo-Schreibweise:
\[
    \beta\tilde{p} + \alpha r \equiv 0 \pmod D
\]
Wir lösen das nach $\tilde{p}$ auf: $\beta\tilde{p} \equiv -\alpha r \pmod D$.
Um das $\beta$ wegzubekommen, multiplizieren wir mit dem modularen Inversen $\beta^{-1}$ (dessen Existenz wir in Schritt 3 beweisen):
\[
    \tilde{p} \equiv -\alpha\beta^{-1} r \pmod D
\]
\textbf{Was bedeutet das?} Auf einer gegebenen Spur $r$ finden wir nicht an jeder beliebigen Stelle Gitterpunkte, sondern nur an den Stellen $\tilde{p}$, die genau in dieses bestimmte Raster (diesen Rest $c_r$ modulo $D$) fallen. Jede Spur liefert also eine Kette von Punkten mit dem konstanten Abstand $D$, beginnend bei einem Offset $c_r$.

\section{Schritt 3: Teilerfremdheit beweisen}

Durften wir oben einfach $\beta^{-1}$ benutzen? Ja, wenn $\gcd(\beta, D) = 1$. 
Wir wissen $D = \alpha^2 + \beta^2$. Wenn $\alpha$ und $\beta$ teilerfremd sind (Grundvoraussetzung), dann kann $\beta$ keinen gemeinsamen Teiler mit $\alpha^2 + \beta^2$ haben (außer 1). Daher existiert $\beta^{-1} \pmod D$ immer.

\section{Schritt 4: Die Verteilung der Spuren auf dem Zahnrad}

Wir haben $N$ Spuren (also $N$ fortlaufende Werte für $r$). Jede Spur erzeugt Punkte auf der Achse, die um $c_r = -\alpha\beta^{-1} r \pmod D$ verschoben sind.
Wir nennen die Konstante $m = -\alpha\beta^{-1}$. 

Wenn $r$ um 1 hochzählt (wir gehen zur nächsten Spur), springt der Offset $c_r$ auf unserem $D$-Zahnrad um genau $m$ Zähne weiter. 
Da auch $m$ teilerfremd zu $D$ ist, greift wieder die Getriebe-Analogie: Wenn wir Spur für Spur durchgehen, überspringen wir zwar scheinbar zufällig Zähne, aber wir treffen \textbf{jeden Zahn genau einmal}, bevor wir wieder am Anfang landen! Die Zuordnung ist eine Bijektion (eine 1:1-Zuordnung).

Wir machen das für $N$ Spuren. 
Wenn wir $N$ durch $D$ teilen ($N = q \cdot D + s$), drehen wir uns $q$ Mal komplett um das Zahnrad und machen noch $s$ Restschritte.
\begin{itemize}
    \item $s$ Zähne auf dem Rad wurden $(q+1)$ mal getroffen.
    \item Die restlichen $(D-s)$ Zähne wurden nur $q$ mal getroffen.
\end{itemize}
Das Wichtige: Da wir in der Abfolge $r, r+1, r+2, \dots$ vorgehen, bilden die $s$ "schweren" Zähne (die einen Treffer mehr haben) einen \textbf{zusammenhängenden Block} (ein \textit{cyclic interval}) in der Sprungabfolge des Zahnrads!

\section{Schritt 5: Beweis der Minimalität}

Jetzt kommt der brillante logische Schluss des Papers (Widerspruchsbeweis):

Angenommen, die Wiederholungs-Periode der Abstände wäre nicht $N$, sondern eine kleinere Zahl $\lambda' < N$. Da $N$ definitiv eine Periode ist, müsste $\lambda'$ ein Teiler von $N$ sein (z.B. $N=10$, und wir vermuten die Periode ist 5).

Wenn das Muster wirklich die Periode $\lambda'$ hat, dann summiert sich die Länge von $\lambda'$ aufeinanderfolgenden Gaps (Abständen) zu einer festen Verschiebungs-Distanz $\sigma$. Das würde bedeuten: Wenn wir unser gesamtes Punktemuster um exakt $\sigma$ verschieben, deckt es sich wieder zu 100\% mit sich selbst.

Aber wenn wir das Muster verschieben, verschieben wir auch unser "Zahnrad" der Offsets um $\sigma$.
Damit das Muster identisch bleibt, muss das Zahnrad nach der Verschiebung exakt gleich aussehen. Die "schweren" Zähne müssen exakt wieder auf "schweren" Zähnen landen.

Wir haben aber in Schritt 4 festgestellt, dass die schweren Zähne einen \textbf{zusammenhängenden Block} bilden. Es gibt keine Symmetrie in einem einzelnen zusammenhängenden Block auf einem Kreis! Wenn Sie auf einem Kuchennetz ein durchgehendes Stück von z.B. 3 von 8 Stücken markieren, können Sie den Kuchen nicht verdrehen, sodass die 3 Stücke wieder genau auf sich selbst liegen (außer Sie drehen um $0^\circ$ oder volle $360^\circ$).

Die einzig mögliche Verschiebung $\sigma$, die den Block auf sich selbst abbildet, ist ein Vielfaches von $D$ (eine volle Umdrehung). Das entspricht aber genau unserer Verschiebung um $N$ Punkte. 

\textbf{Folgerung:} Es kann keine kleinere Periode $\lambda'$ geben, die das Muster in sich selbst überführt. Die Periode \textbf{muss} minimal $N$ sein. (Q.E.D.)

\end{document}
